\documentclass[12pt,a4paper]{article}

\usepackage[T1]{fontenc}
\usepackage[utf8]{inputenc}
\usepackage[brazil]{babel}
\usepackage{lmodern}
\usepackage{geometry}
\usepackage{microtype}
\usepackage{booktabs}
\usepackage{array}
\usepackage{graphicx}
\usepackage{float}
\usepackage{xcolor}
\usepackage{listings}
\usepackage{hyperref}
\usepackage{amsmath}
\usepackage{enumitem}

\geometry{margin=2.5cm}
\hypersetup{colorlinks=true,linkcolor=blue!45!black,urlcolor=blue!55!black}
\setlength{\parindent}{1.25cm}
\setlength{\parskip}{0.25em}
\renewcommand{\arraystretch}{1.2}

% Preencha apenas os dados de identificação ainda ausentes.
\newcommand{\NomeAluno}{Luiz Alberto Silva Mota}
\newcommand{\RAAluno}{PREENCHER RA}
\newcommand{\CursoAluno}{PREENCHER CURSO}
\newcommand{\DataEntrega}{22 de setembro de 2026}

\newcommand{\figuraouespaco}[3]{%
  \begin{figure}[H]
    \centering
    \IfFileExists{#1}{\includegraphics[width=#2\textwidth]{#1}}{%
      \fbox{\parbox[c][5.2cm][c]{0.88\textwidth}{\centering
      Espaço reservado para #3\\[0.5em]\texttt{#1}}}}
    \caption{#3}
  \end{figure}}

\lstset{
  language=Python,
  basicstyle=\ttfamily\small,
  keywordstyle=\color{blue!60!black},
  commentstyle=\color{green!40!black},
  stringstyle=\color{orange!60!black},
  frame=single,
  breaklines=true,
  showstringspaces=false
}

\begin{document}

\begin{titlepage}
  \centering
  {\large Universidade Presbiteriana Mackenzie\par}
  \vspace{1.4cm}
  {\large \CursoAluno\par}
  {\large Disciplina: Computação Distribuída\par}
  \vfill
  {\LARGE\bfseries Processamento Distribuído de Imagens Médicas com MPI\par}
  \vspace{0.8cm}
  {\large Relatório do Laboratório 04\par}
  \vfill
  \begin{tabular}{rl}
    Aluno: & \NomeAluno \\
    RA: & \RAAluno \\
    Docentes: & Prof. Alcides / Prof. Mário \\
  \end{tabular}
  \vfill
  São Paulo\\
  \DataEntrega
\end{titlepage}

\tableofcontents
\newpage

\section{Introdução e fundamentos teóricos}

Radiografias digitais em tons de cinza podem ser representadas por matrizes
bidimensionais, nas quais cada elemento contém a intensidade de um pixel no
intervalo de 0 a 255. Em uma triagem computacional simplificada, intensidades
elevadas no interior dos campos pulmonares são tratadas como possíveis regiões
de consolidação. O objetivo deste trabalho é detectar e localizar essas
alterações em uma radiografia sintética, reduzindo o tempo de processamento por
paralelismo de dados.

O padrão MPI organiza processos independentes em um comunicador. Cada processo
possui um identificador (\emph{rank}) e o processo de rank zero atua como
coordenador. A imagem é dividida em faixas horizontais contíguas e cada rank
analisa somente sua faixa. A solução utiliza cinco operações centrais:
\texttt{Bcast} difunde parâmetros; \texttt{Scatter} distribui as faixas;
\texttt{Barrier} delimita fases; \texttt{Reduce} agrega métricas numéricas; e
\texttt{Gather} coleta relatórios descritivos individuais.

O modelo é SPMD (\emph{Single Program, Multiple Data}): todos executam o mesmo
programa, mas operam sobre dados diferentes. A aceleração é limitada pelas
partes seriais, pelo custo de comunicação e pelo rank mais lento, conforme a
Lei de Amdahl e a semântica das sincronizações coletivas.

\section{Metodologia e arquitetura do cluster}

\subsection{Infraestrutura}

O ambiente é reproduzido com Docker Compose e contém quatro contêineres: um
\texttt{master} e três \texttt{workers}. Cada contêiner anuncia dois slots no
arquivo \texttt{hosts}, totalizando oito processos. A imagem base é Python 3.12
sobre Debian Bookworm, com Open MPI, OpenSSH, \texttt{mpi4py} 4.1.0, NumPy 2.2.6
e Matplotlib 3.10.3. O mesmo diretório do projeto é montado em todos os nós, e o
Open MPI inicia processos remotos por SSH sem senha em uma rede Docker isolada.

\subsection{Geração da radiografia sintética}

Somente o rank zero cria a matriz. O fundo e o mediastino recebem intensidades
baixas/médias com gradiente crânio-caudal. Duas máscaras elípticas representam
os campos pulmonares e recebem valores predominantemente normais. Uma fração
configurável dos pixels pulmonares é então selecionada, sem reposição, para
receber intensidades a partir do limiar suspeito. A semente pseudoaleatória
fixa em 2026 garante que a imagem seja idêntica nas comparações entre números de
processos.

Os valores padrão são: limiar suspeito 200, limiar alto 230, percentual crítico
5\% e taxa sintética de anomalia pulmonar 4\%. Uma faixa é classificada como
\texttt{NORMAL} abaixo de 1\% de pixels suspeitos; \texttt{ATENCAO} entre 1\% e
5\%; e \texttt{CRITICA} a partir de 5\%. Para tornar objetiva a expressão
``presença expressiva de pixels altamente suspeitos'', também se considera
crítica a faixa em que pixels acima de 230 atinjam metade do percentual crítico.

\subsection{Fluxo distribuído}

O rank zero empacota dimensões, limiares, semente e parâmetros de atraso em um
dicionário e executa \texttt{comm.bcast}. Uma primeira barreira garante que
todos concluíram a parametrização. A matriz é dividida por
\texttt{numpy.array\_split}, que preserva todas as linhas mesmo quando a divisão
não é exata, e os blocos são distribuídos por \texttt{comm.scatter}.

Cada rank calcula quantidade e soma de pixels, máximo local, contagens suspeita
e altamente suspeita e contagens por pulmão esquerdo e direito. Para simular
heterogeneidade, ranks ímpares aguardam $0{,}5\times rank$ segundos após o
processamento local; ranks pares não recebem atraso. A segunda barreira impede
que a consolidação comece antes do rank mais lento.

As somas e o máximo globais são obtidos com \texttt{comm.reduce}, usando
\texttt{MPI.SUM} e \texttt{MPI.MAX}. Em paralelo, cada rank monta um dicionário
com intervalo de linhas, dimensão, média, pico, classificação e tempo local;
esses registros são enviados ao coordenador por \texttt{comm.gather}. Por fim,
o rank zero calcula percentuais, lateralidade e diagnóstico geral, imprime o
relatório e opcionalmente grava JSON.

\section{Evidências experimentais}

Os experimentos devem usar a mesma semente e parâmetros em todas as execuções.
As capturas abaixo comprovam o funcionamento com as três quantidades de
processos exigidas.

\figuraouespaco{figuras/execucao_2_processos.png}{0.95}{Execução com 2 processos MPI}
\figuraouespaco{figuras/execucao_4_processos.png}{0.95}{Execução com 4 processos MPI}
\figuraouespaco{figuras/execucao_8_processos.png}{0.95}{Execução com 8 processos MPI}

As três execuções produziram o mesmo resultado clínico para a imagem de
$2000\times2000$: 58.754 pixels suspeitos (1,4688\%), 27.118 altamente
suspeitos (0,6780\%), maior intensidade igual a 255 e classificação geral
\texttt{ATENCAO}. O pulmão esquerdo apresentou 29.472 pixels suspeitos, contra
29.282 no direito. A Figura 3 contém também o relatório consolidado completo e
a auditoria dos oito ranks.

\section{Desempenho e escalabilidade}

Cada configuração foi executada três vezes. A Tabela 1 apresenta a média dos
tempos totais registrados no campo \texttt{tempo\_total\_ms} dos JSONs. A linha
de um processo é a referência para calcular o speedup clássico.

\begin{table}[H]
\centering
\caption{Tempo médio de execução em milissegundos.}
\begin{tabular}{ccccc}
\toprule
Dimensão & 1 processo & 2 processos & 4 processos & 8 processos \\
\midrule
$2000\times2000$ & $740{,}602$ & $819{,}135$ & $1.972{,}004$ & $4.111{,}979$ \\
$4000\times4000$ & $3.691{,}355$ & $1.832{,}817$ & $3.004{,}286$ & $6.393{,}504$ \\
$6000\times6000$ & $6.996{,}085$ & $2.659{,}646$ & $4.584{,}165$ & $8.230{,}646$ \\
\bottomrule
\end{tabular}
\end{table}

O speedup e a eficiência são definidos por
\begin{equation}
S_p=\frac{T_1}{T_p}, \qquad E_p=\frac{S_p}{p}.
\end{equation}
O CSV registrou speedups de 0,904, 0,376 e 0,180 para a imagem
$2000\times2000$ com 2, 4 e 8 processos. Para $4000\times4000$, os valores
foram 2,014, 1,229 e 0,577; para $6000\times6000$, 2,630, 1,526 e 0,850.

\figuraouespaco{figuras/speedup.png}{0.80}{Speedup por dimensão da imagem}

A comparação deve considerar conjuntamente tempo, speedup e eficiência. Para
$2000\times2000$, um processo apresentou a menor média. Para as matrizes
$4000\times4000$ e $6000\times6000$, dois processos apresentaram as menores
médias, respectivamente 1.832,817 ms e 2.659,646 ms. Quatro e oito processos
foram prejudicados pelos atrasos dos ranks ímpares, especialmente pelo rank 7,
que adiciona 3,5 s. A dispersão do baseline de um processo foi elevada:
desvios-padrão de 659,979 ms, 4.095,725 ms e 7.751,170 ms. Por isso, os speedups
superlineares de 2,014 e 2,630 devem ser tratados com cautela, pois refletem
também instabilidade ou efeito de aquecimento nas execuções sequenciais.

\section{Questões de análise e reflexão}

\subsection*{1. Impacto do número de processos no tempo de execução}
O tempo não precisa diminuir linearmente. A fração serial inclui geração da
imagem, criação dos blocos, impressão e escrita do JSON. Somam-se a ela
serialização, transferência, criação de processos e sincronização. Pela Lei de
Amdahl, mesmo uma parte serial pequena limita o speedup máximo. Para a imagem
$2000\times2000$, a média aumentou de 819,135 ms com dois processos para
1.972,004 ms com quatro e 4.111,979 ms com oito. Isso representa aumentos de
aproximadamente 140,7\% e 108,5\% nas duas transições, explicados principalmente
pelo atraso artificial dos ranks ímpares e pelas barreiras globais.

\subsection*{2. Influência dos processos lentos (stragglers)}
Depois da análise local, ranks ímpares dormem por $0{,}5\times rank$ segundos.
Como todos chegam à segunda barreira e às operações coletivas, o tempo global é
ditado pelo participante mais lento. Com oito processos, o rank 7 acrescenta
3,5 s; os demais permanecem ociosos na barreira. Isso ilustra que balancear a
quantidade de dados não elimina heterogeneidade de hardware.

\subsection*{3. Papel prático de \texttt{MPI\_Barrier}}
A primeira barreira separa configuração e distribuição; a segunda separa
análise local e consolidação. A segunda é especialmente útil no experimento,
pois torna observável o impacto do straggler e impede avanço de fase. Sem ela,
as próprias operações coletivas ainda imporiam dependências, mas a fronteira
entre fases ficaria implícita e os tempos seriam mais difíceis de interpretar.

\subsection*{4. \texttt{Bcast} versus \texttt{Scatter}}
Os parâmetros são pequenos e todos os ranks precisam da mesma cópia, portanto
\texttt{Bcast} é adequado. A imagem é grande e cada rank precisa apenas de uma
faixa, portanto \texttt{Scatter} evita replicação. Enviar a imagem completa por
\texttt{Bcast} multiplicaria tráfego e memória por processo e exigiria que cada
rank descartasse dados desnecessários.

\subsection*{5. Diferença prática entre \texttt{Reduce} e \texttt{Gather}}
\texttt{Reduce} combina valores por uma operação associativa, sendo apropriado
para soma, contagem e máximo sem preservar a origem de cada parcela.
\texttt{Gather} mantém um registro por rank e permite auditoria das faixas,
classificações e tempos. Ambos são necessários: o primeiro produz métricas
globais eficientes; o segundo conserva a individualidade dos nós.

\subsection*{6. Escalabilidade em imagens pequenas}
Para $500\times500$, a carga útil por rank é pequena. Inicialização do MPI,
SSH, serialização, cópia, barreiras e reduções podem custar mais do que a
varredura NumPy. Assim, distribuir a análise tende a não trazer vantagem e pode
até aumentar o tempo. A bateria registrada começa em $2000\times2000$ e já
mostra essa tendência: oito processos demoraram 4.111,979 ms, contra 740,602 ms
com um processo. Em $500\times500$, a relação entre overhead e trabalho útil é
ainda menos favorável.

\subsection*{7. Detecção de lateralidade anatômica}
As colunas menores que $\lfloor N/2\rfloor$ representam o pulmão esquerdo e as
demais representam o direito. Cada rank recebe faixas horizontais, mas conserva
todas as colunas; portanto a coordenada horizontal permanece válida em todos os
blocos. Essa geometria permite somar separadamente suspeitos dos dois lados e
inferir qual está mais comprometido.

\subsection*{8. Balanceamento de carga}
\texttt{array\_split} distribui números de linhas iguais ou diferentes por no
máximo uma unidade, logo a carga nominal é uniforme. Entretanto, a quantidade
de suspeitos pode variar entre faixas apicais, médias e basais devido à posição
aleatória das alterações; o atraso artificial também cria desequilíbrio de
tempo. Na execução com oito processos, as contagens variaram de 45 pixels
suspeitos no rank 7 a 11.918 no rank 3, uma diferença de 11.873 pixels. Os
tempos locais variaram de 0,741 ms a 3.511,767 ms, uma diferença de
3.511,026 ms determinada principalmente pelo atraso artificial.

\subsection*{9. Desafio de divisibilidade de linhas}
Para 2003 linhas e quatro processos, a divisão correta é 501, 501, 501 e 500
linhas. \texttt{numpy.array\_split} aplica essa distribuição automaticamente,
sem \emph{padding} e sem perda. O início da faixa é calculado por
$rank\lfloor M/P\rfloor+\min(rank,M\bmod P)$, mantendo intervalos contíguos e
sem sobreposição. O teste deve registrar exatamente
$2003\times2003=4\,012\,009$ pixels analisados, confirmando a ausência de perda
ou duplicação.

\section{Dificuldades técnicas e soluções}

A execução MPI em múltiplos contêineres exige resolução dos nomes dos serviços,
servidor SSH ativo e autenticação sem senha. A imagem Docker gera uma chave para
o usuário \texttt{mpiuser}, instala Open MPI em todos os nós e usa uma rede
Compose comum. Outro ponto é a divisibilidade: uma divisão inteira simples
perderia o resto, enquanto \texttt{array\_split} mantém todas as linhas. Por
fim, valores NumPy foram convertidos em inteiros Python antes de reduções e do
JSON, evitando incompatibilidades de serialização. A automação também verifica
o código de saída de cada execução, preserva o terminal em arquivos de log e
só então gera o resumo, evitando incorporar medições incompletas.

\section{Declaração e análise do uso de IA}

Foi utilizada a ferramenta OpenAI Codex/ChatGPT principalmente como apoio na
organização e redação deste relatório e no esclarecimento das questões de
análise e reflexão que apresentaram maior dificuldade. O suporte ajudou a
estruturar as respostas, relacionar os resultados experimentais aos conceitos
de comunicação coletiva, sincronização, balanceamento de carga, speedup e Lei
de Amdahl, além de revisar a apresentação dos dados obtidos nos testes.

Também houve apoio na revisão da solução técnica e do roteiro de execução. As
sugestões foram verificadas a partir do código, dos arquivos JSON e CSV e das
capturas do terminal. Foram conferidos o uso de \texttt{Bcast},
\texttt{Scatter}, \texttt{Barrier}, \texttt{Reduce} e \texttt{Gather}, a
divisão das linhas com \texttt{array\_split} e a coerência entre as métricas e
os resultados apresentados. Dessa forma, as respostas não foram aceitas
automaticamente: foram comparadas com a implementação e com as evidências
experimentais antes de serem incorporadas ao relatório. O principal aprendizado
foi compreender a relação entre os coletivos MPI, os custos de sincronização e
o comportamento de desempenho observado.

\section{Conclusão}

A solução decompõe uma radiografia sintética em faixas horizontais, mantém a
semântica anatômica das colunas e consolida um diagnóstico auditável. A
implementação cobre comunicação, sincronização, redução e coleta, além de lidar
com divisões não exatas e heterogeneidade. Os resultados mostram que aumentar
processos só é vantajoso quando a carga útil compensa o overhead e quando nenhum
straggler domina a fase coletiva. Um processo foi melhor em
$2000\times2000$, enquanto dois processos foram melhores em $4000\times4000$ e
$6000\times6000$. O resultado confirma que paralelismo adicional só compensa
quando o trabalho útil supera comunicação, sincronização e heterogeneidade.

\begin{thebibliography}{9}
\bibitem{mpi}
MPI Forum. \emph{MPI: A Message-Passing Interface Standard, Version 4.1}.
Disponível em: \url{https://www.mpi-forum.org/docs/mpi-4.1/mpi41-report.pdf}.

\bibitem{mpi4py}
DALCÍN, Lisandro. \emph{MPI for Python Documentation}.
Disponível em: \url{https://mpi4py.readthedocs.io/}.

\bibitem{numpy}
NUMPY DEVELOPERS. \emph{NumPy Reference: array\_split}.
Disponível em: \url{https://numpy.org/doc/stable/reference/generated/numpy.array_split.html}.

\bibitem{docker}
DOCKER. \emph{Docker Compose Documentation}.
Disponível em: \url{https://docs.docker.com/compose/}.
\end{thebibliography}

\end{document}
